    \section{Autores}

    As vezes, dizer obrigado é importante, esse livro tem o nome e o coração de um grupo, cheio de discordâncias. Nós somos \textbf{diferentes} e isso é o que nos faz \textbf{fortes}. Acredite em si mesmo, e se você não consegue, acredite nas pessoas ao seu lado.

    Obrigado.

    Aqui uma lista de nomes que contribuíram com esse livro (para aqueles que estão se perguntando, isso foi ordenado usando random\_shuffle no C++20 usando srand(252)):
    \begin{itemize}
        \item Lesin
        \item Alunea
        \item Nopebi Lifesa
        \item Atak Kichan
        \item Faslecar
        \item Nollyad
        \item Laelovatsug
        \item Ognol Tohberum
        \item Nhotivew
    \end{itemize}
    \section{Dicas}
    \begin{itemize}
    \item Sempre procure por invariantes no problema, a soma do vetor muda? O que acontece com o número de inversões? Esse tipo de coisa.

    \item Se o problema parece impossível, então talvez você tenha que transformar o problema em outra coisa pra aí ficar possível.

    \item Se você quer dividir o seu universo em dois, ao invés de tentar dividir os dois ao mesmo tempo, imagina que você vai colocar todo mundo junto e daí você vai tirar uns caras do todo.
    
    \item Lembre-se que binary lifting não é apenas para árvores.
    
    \item Valor esperado? Contribuição é o caminho.
    
    \item DP? Talvez uma solução lenta possa ser otimizada.
    
    \item Relaxa, a competição é longa, começar devagar é bom...
    
    \item Respire, beba água, faça piadas, coma doces, no final, se divirta com seus amigos.
    
    \item Um guloso é arriscado, mas dá pra fazer. Ninguém sabe como provar essa bomba mesmo...
    
    \item Fluxo? É bipartido? Lembre-se do que dá para fazer em grafos bipartidos:
    \begin{itemize}
        \item Cobertura mínima de vértices = Matching máximo;
        \item Clique máxima = Matching máximo.
        \item Conjunto independente máximo = Complemento da cobertura mínima de vértices.
    \end{itemize}
    
    \item Lembre-se, uma busca binária pode ajudar em muito!

    \item Consegue modelar alguma recorrência?

    \item Escreva muito, Pense para fora, Escreva.

    \item Não fique bitolado em achar uma solução rápida, ache uma solução e, então, tente entendê-la. Tente comer as beiradas, desenhe casos e entenda o problema.

    \item Tem módulo? Talvez dividir em casos ajude.

    \item É monotônico (crescente ou decrescente)?
    
    \item Geometria tem 4 opções: Line Sweep, busca binária, Convex Hull e matemática.

    \item Teoria dos jogos? Apenas brute. Recorrência em números grandes? Apenas brute. Brutar é god d+.
    
    \item RESPIRE!
    
    \item Seus amigos estão aqui para ajudar!
    
    \item Abraçe uns aos outros, espalhe amor, não guerra.
    
    \item Uma piada não tem que ofender um amigo.
    
    \item VOCÊS SÃO UM TIME!
    
    \item Talvez checar apenas os primos é o suficiente (se não funcionar, tente potências de dois).

    \item Se você não consegue fazer em $n\log(n)$, apenas faz $n \log(n)^2$ seu bobo.
    
    \item Lembre-se, você pode fazer uns fatoriais usando convoluções.

    \item Sem tunnel vision, se você tá empacado, dê tempo para si mesmo.

    \item WA? Reflita a sua ideia, depois leia seu código INTEIRO e então bruta isso (PELO AMOR DE DEUS, BRUTA).
    
    \end{itemize}
